% source: Allen B. Altman and Steven L. Kleiman, "Compactifying the Picard Scheme", Adv. Math. 35 (1980)
% pdf_page: 0018
% doc_page: 68
% retrieved: 2026-06-17
% transcribed_by: codex/gpt-5.6-sol (vision, rendered at 200dpi via pdftoppm; no OCR)

% requires: amsmath

\DeclareMathOperator{\coker}{coker}

% running head: ALTMAN AND KLEIMAN / p. 68
% [continuation of the proof of Theorem (2.6)]

\paragraph{Step III.}
Let $g:T\to\mathcal{G}$ be an $S$-morphism, and let $G$ be an element of
$\mathbf{A}(T)$. Set $G'=(1\times g)^*H$. Then $G'$ is equivalent to $G$ as
a quotient of $F_T$ if and only if the $\phi(m)$-quotients $g^*Q$ and
$(f_T)_*G(m)$ of $(f_*F(m))_T$ are equivalent.

\paragraph{Proof.}
Suppose $g^*Q$ and $(f_T)_*G(m)$ are equivalent $\phi(m)$-quotients. Then
there is a diagram with exact rows and commutative right-hand square,
\[
\begin{array}{ccccccccc}
0&\longrightarrow&g^*K&\xrightarrow{g^*(u)}&
g^*(f_{\mathcal G})_*F_{\mathcal G}(m)&\longrightarrow&g^*Q&\longrightarrow&0
\\
&&\underset{\sim}{\big\downarrow}\,\searrow^{v}
&&\underset{\sim}{\big\downarrow}
&&\underset{\sim}{\big\downarrow}
\\
0&\longrightarrow&(f_T)_*I(m)&\longrightarrow&(f_T)_*F_T(m)
&\longrightarrow&(f_T)_*G(m)&\longrightarrow&0,
\end{array}
\]
where $I$ is the pseudo-ideal of $G$ and the middle map is the base-change
isomorphism. The bottom row is exact because, since $I$ is $m$-regular on
the fibers, $R^1(f_T)_*I(m)$ is equal to zero. Hence the dotted isomorphism
making the left-hand square commutative exists.

Taking the adjoint of the lower left-hand triangle yields the commutative
diagram,
\[
\begin{array}{ccc}
\multicolumn{3}{c}{(f_T)^*g^*K=(1\times g)^*f_{\mathcal G}^*K}
\\
\underset{\sim}{\big\downarrow}&\searrow^{v^\#}&
{\scriptstyle (u^\#)_T}\big\downarrow
\\
(f_T)^*(f_T)_*I(m)&\longrightarrow&F_T(m).
\end{array}
\]
Since $I$ is $m$-regular on the fibers, the canonical map,
\[
(f_T)^*(f_T)_*I(m)\longrightarrow I(m),
\]
is surjective by base-change theory and by [SGA 6, XIII, 1.3(iii), p. 616].
So the image of the lower horizontal map is equal to $I(m)$. Hence the
quotient of $F_T(n)$ it defines is $G(m)$. On the other hand, $G'(m)$ is
clearly equal to $\coker((u^\#)_T)$. Thus $G'$ is equivalent to $G$.
% [sic: the printed text has F_T(n) in the preceding sentence]

For the converse, start with the diagram with exact rows and commutative
right square,
\[
\begin{array}{ccccccccc}
&& (f_{\mathcal G}^*K)_T&\xrightarrow{(u^\#)_T}&F_T(m)
&\longrightarrow&G'(m)&\longrightarrow&0
\\
&&{\scriptstyle\mathrm{dotted}}\big\downarrow\,\searrow
&&\big\downarrow{\scriptstyle =}
&&\big\downarrow
\\
0&\longrightarrow&I(m)&\longrightarrow&F_T(m)
&\longrightarrow&G(m)&\longrightarrow&0.
\end{array}
\]
Induced is the dotted vertical map making the left-hand square commutative.
Taking the adjoint of the lower left-hand triangle yields the diagram with
exact rows and commutative left square,
% [the diagram continues on page-0019]
